% F085｜真实数学图；代表函数不构成抽象解的唯一性证明。
\begin{center}
\begin{minipage}{.49\linewidth}\centering
\begin{tikzpicture}
\begin{axis}[width=6.1cm,height=4.8cm,axis lines=middle,
 xmin=-3.7,xmax=3.7,ymin=-1.4,ymax=1.6,
 xtick={-3,0,3},xticklabels=\empty,ytick={-1,0,1},xlabel={$t$},ylabel={$f(t)$},
 title={余弦代表：关于纵轴对称},title style={font=\small,yshift=3pt},
 tick label style={font=\small},label style={font=\small},
 ylabel style={at={(axis description cs:0,1)},anchor=south east,xshift=-18pt},
 clip=false,grid=major,grid style={black!8},enlargelimits=false,
 every axis plot/.append style={thick},samples=161]
\addplot[CExplanation,very thick,domain=-3.14:3.14] {cos(deg(x))};
\draw[dashed,black!35] (axis cs:-1.0471975511965976,0) |- (axis cs:0,0.5);
\draw[dashed,black!35] (axis cs:1.0471975511965976,0) |- (axis cs:0,0.5);
\addplot[only marks,mark=*,mark size=2pt,orange!80!black] coordinates {(-1.04719755,0.5) (1.04719755,0.5)};

\node[font=\small,anchor=north,yshift=-3pt] at (axis cs:-3,-1.4) {$-3$};
\node[font=\small,anchor=north,yshift=-3pt] at (axis cs:3,-1.4) {$3$};

\end{axis}
\end{tikzpicture}
\end{minipage}\hfill
\begin{minipage}{.49\linewidth}\centering
\begin{tikzpicture}
\begin{axis}[width=6.1cm,height=4.8cm,axis lines=middle,
 xmin=-3.7,xmax=3.7,ymin=-1.5,ymax=1.5,
 xtick={-3,0,3},xticklabels=\empty,ytick={-1,0,1},xlabel={$t$},ylabel={$f(t)$},
 title={正弦代表：关于原点对称},title style={font=\small,yshift=3pt},
 tick label style={font=\small},label style={font=\small},
 ylabel style={at={(axis description cs:0,1)},anchor=south east,xshift=-18pt},
 clip=false,grid=major,grid style={black!8},enlargelimits=false,
 every axis plot/.append style={thick},samples=161]
\addplot[CExplanation,very thick,domain=-3.14:3.14] {sin(deg(x))};
\draw[dashed,black!35] (axis cs:-0.5235987755982988,0) |- (axis cs:0,-0.5);
\draw[dashed,black!35] (axis cs:0.5235987755982988,0) |- (axis cs:0,0.5);
\addplot[only marks,mark=*,mark size=2pt,orange!80!black] coordinates {(-0.523598776,-0.5) (0.523598776,0.5)};

\node[font=\small,anchor=north,yshift=-3pt] at (axis cs:-3,-1.5) {$-3$};
\node[font=\small,anchor=north,yshift=-3pt] at (axis cs:3,-1.5) {$3$};

\end{axis}
\end{tikzpicture}
\end{minipage}
\par\small 特殊值代入得到的是对称关系。这里用代表函数展示，证明对抽象模型直接进行。
\end{center}
